\chapter{Exceptions Système}
\sloppy

\section{Exercice 1 -- Division Sécurisée}

\subsection{Objectif}

Calculer la division de deux entiers saisis par l'utilisateur en gérant
la division par zéro et la saisie invalide.

\subsection{Code source}

\begin{lstlisting}[language=Java, caption={exo1.java -- Division sécurisée}]
import java.util.Scanner;
import java.util.InputMismatchException;

public class exo1 {

    public exo1() {}

    public void calculer(exo1 e, int a, int b) {
        try {
            if (b == 0) {
                throw new ArithmeticException();
            }
            System.out.println(a / b);
        } catch (ArithmeticException ex) {
            System.out.println("Erreur : Division par zero.");
        } catch (Exception ex) {
            System.out.println("Erreur inattendue.");
        }
    }

    public static void main(String[] args) {
        Scanner sc = new Scanner(System.in);
        try {
            System.out.print("Enter a : ");
            int a = sc.nextInt();
            System.out.print("Enter b : ");
            int b = sc.nextInt();
            exo1 d = new exo1();
            d.calculer(d, a, b);
        } catch (InputMismatchException ex) {
            System.out.println("Erreur : Saisie invalide, entrez des entiers.");
        }
        sc.close();
    }
}
\end{lstlisting}

\subsection{Explication}

Le bloc \texttt{try/catch} dans la méthode \texttt{calculer()} intercepte
l'\texttt{ArithmeticException} levée explicitement via \texttt{throw} lorsque
le diviseur est nul. Dans le \texttt{main}, un second \texttt{try/catch}
capture l'\texttt{InputMismatchException} levée par \texttt{Scanner} en cas de
saisie non entière.

\begin{infobox}
    \centering\vspace{0.2cm}
    {\large Le mot-clé \texttt{throw} permet de lever manuellement une exception
    depuis n'importe quel point du code.}
    \vspace{0.2cm}
\end{infobox}

\section{Exercice 2 -- NullPointerException}

\subsection{Objectif}

Afficher la longueur d'une chaîne en gérant le cas où elle est \texttt{null},
d'abord sans \texttt{try/catch} puis avec.

\subsection{Code source}

\begin{lstlisting}[language=Java, caption={exo2.java -- NullPointerException}]
public class exo2 {

    public void longueurSansCatch(exo2 e, String s) {
        if (s == null) {
            System.out.println("Null");
        } else {
            System.out.println(s.length());
        }
    }

    public void longueurAvecCatch(exo2 e, String s) {
        try {
            System.out.println(s.length());
        } catch (NullPointerException ex) {
            System.out.println("Null");
        }
    }
}
\end{lstlisting}

\subsection{Explication}

La version sans \texttt{try/catch} utilise une condition \texttt{if} pour
tester la nullité avant d'appeler la méthode. La version avec \texttt{try/catch}
laisse la \texttt{NullPointerException} se lever et la capture dans le bloc
\texttt{catch}.

\section{Exercice 3 -- Tableau Sécurisé}

\subsection{Objectif}

Permettre à l'utilisateur d'accéder à un index d'un tableau de taille 5
en proposant deux approches.

\subsection{Code source}

\begin{lstlisting}[language=Java, caption={exo3.java -- Accès sécurisé au tableau}]
public class exo3 {

    public void versionIf(exo3 e, int[] tab, int index) {
        if (index >= 0 && index < 5) {
            System.out.println(tab[index]);
        } else {
            System.out.println("Hors limites");
        }
    }

    public void versionCatch(exo3 e, int[] tab, int index) {
        try {
            System.out.println(tab[index]);
        } catch (ArrayIndexOutOfBoundsException ex) {
            System.out.println("Hors limites");
        }
    }
}
\end{lstlisting}

\subsection{Comparaison des deux approches}

\begin{table}[H]
\centering
\begin{tabularx}{0.85\textwidth}{C{3cm} X X}
\toprule
\textbf{\color{primarycolor}Critère} &
\textbf{\color{primarycolor}Version if} &
\textbf{\color{primarycolor}Version try/catch} \\
\midrule
Lisibilité & Explicite & Nécessite de connaître l'exception \\
\addlinespace[0.15cm]
Performance & Légèrement plus rapide & Surcoût si exception levée \\
\addlinespace[0.15cm]
Usage recommandé & Vérification préventive & Cas imprévisibles \\
\bottomrule
\end{tabularx}
\caption{Comparaison entre vérification par \texttt{if} et \texttt{try/catch}}
\end{table}

\section{Exercice 4 -- Conversion d'Entier}

\subsection{Objectif}

Convertir une chaîne saisie en entier et afficher un message clair en cas d'échec.

\subsection{Code source}

\begin{lstlisting}[language=Java, caption={exo4.java -- NumberFormatException}]
public class exo4 {

    public void convertir(exo4 e, String s) {
        try {
            int val = Integer.parseInt(s);
            System.out.println(val);
        } catch (NumberFormatException ex) {
            System.out.println("Erreur : La chaine saisie ne peut pas etre convertie en entier.");
        }
    }
}
\end{lstlisting}

\subsection{Explication}

\texttt{Integer.parseInt()} lève une \texttt{NumberFormatException} si la
chaîne ne représente pas un entier valide. Ce type d'exception est une
\textit{unchecked exception} (hérite de \texttt{RuntimeException}).

\section{Exercice 5 -- Choix d'Exception}

\subsection{Objectif}

Créer une méthode \texttt{racineCarree(int x)} qui lève une exception si \texttt{x < 0}.

\subsection{Code source}

\begin{lstlisting}[language=Java, caption={exo5.java -- IllegalArgumentException}]
public class exo5 {

    public int racineCarree(exo5 e, int x) {
        if (x < 0) {
            throw new IllegalArgumentException();
        }
        return (int) Math.sqrt(x);
    }
}
\end{lstlisting}

\subsection{Justification du choix}

\begin{accentbox}
    \centering\vspace{0.3cm}
    {\large \texttt{IllegalArgumentException} est choisie car l'erreur provient
    d'un \textbf{argument invalide} passé par le développeur.
    C'est une \textit{unchecked exception} qui ne nécessite pas de
    \texttt{throws} dans la signature, ce qui est adapté à ce type de
    précondition logique.}
    \vspace{0.3cm}
\end{accentbox}

\section{Exercice 6 -- IllegalStateException}

\subsection{Objectif}

Simuler une machine avec un état ON/OFF~; lever une exception si on tente de
la démarrer alors qu'elle est déjà allumée.

\subsection{Code source}

\begin{lstlisting}[language=Java, caption={exo6.java -- IllegalStateException}]
public class exo6 {

    boolean estAllumee;

    public exo6() {
        this.estAllumee = false;
    }

    public void demarrer(exo6 e) {
        if (e.estAllumee) {
            throw new IllegalStateException();
        }
        e.estAllumee = true;
        System.out.println("Machine ON");
    }
}
\end{lstlisting}

\subsection{Explication}

\texttt{IllegalStateException} est l'exception standard Java à utiliser
lorsqu'une méthode est appelée à un moment où l'\textbf{état de l'objet}
ne le permet pas.

\section{Exercice 7 -- Propagation d'Exception}

\subsection{Objectif}

Illustrer la propagation d'exceptions à travers une chaîne d'appels
A $\rightarrow$ B $\rightarrow$ C.

\subsection{Code source}

\begin{lstlisting}[language=Java, caption={exo7.java -- Propagation}]
public class exo7 {

    public void methodeC(exo7 e) throws Exception {
        throw new Exception();
    }

    public void methodeB(exo7 e) throws Exception {
        e.methodeC(e);
    }

    public void methodeA(exo7 e) {
        try {
            e.methodeB(e);
        } catch (Exception ex) {
            System.out.println("Erreur interceptee par A");
        }
    }
}
\end{lstlisting}

\subsection{Schéma de propagation}

\begin{infobox}
    \centering\vspace{0.2cm}
    {\large \texttt{C} lance l'exception $\rightarrow$ \texttt{B} la propage
    (\texttt{throws}) $\rightarrow$ \texttt{A} la capture (\texttt{catch})}
    \vspace{0.2cm}
\end{infobox}

\section{Exercice 8 -- Checked Exception}

\subsection{Objectif}

Créer une méthode \texttt{verifierAge(int age) throws Exception} qui lève une
exception si l'âge est inférieur à 18.

\subsection{Code source}

\begin{lstlisting}[language=Java, caption={exo8.java -- Checked Exception}]
public class exo8 {

    public void verifierAge(exo8 e, int age) throws Exception {
        if (age < 18) {
            throw new Exception();
        }
    }

    public static void main(String[] args) {
        exo8 e = new exo8();
        try {
            e.verifierAge(e, 15);
            System.out.println("Age valide.");
        } catch (Exception ex) {
            System.out.println("Erreur : Age inferieur a 18.");
        }
    }
}
\end{lstlisting}

\subsection{Explication}

En utilisant \texttt{throws Exception}, le compilateur \textbf{oblige}
l'appelant à gérer l'exception. C'est le comportement des
\textit{checked exceptions}.

\section{Exercice 9 -- Exception Personnalisée Vérifiée}

\subsection{Objectif}

Créer une \texttt{AgeException} personnalisée et l'utiliser pour vérifier un âge.

\subsection{Code source}

\begin{lstlisting}[language=Java, caption={exo9.java -- AgeException personnalisée}]
class AgeException extends Exception {}

public class exo9 {

    public void verifierAge(exo9 e, int age) throws AgeException {
        if (age < 18) {
            throw new AgeException();
        }
    }

    public static void main(String[] args) {
        exo9 e = new exo9();
        try {
            e.verifierAge(e, 15);
            System.out.println("Age valide.");
        } catch (AgeException ex) {
            System.out.println("Erreur : Age inferieur a 18.");
        }
    }
}
\end{lstlisting}

\section{Exercice 10 -- Comparaison Checked / Unchecked}

\subsection{Objectif}

Comparer les deux approches~: \texttt{Exception} (checked) et \texttt{RuntimeException} (unchecked).

\subsection{Code source}

\begin{lstlisting}[language=Java, caption={exo10.java -- Comparaison Checked vs Unchecked}]
public class exo10 {

    public void verifierChecked(exo10 e, int age) throws Exception {
        if (age < 18) {
            throw new Exception();
        }
    }

    public void verifierUnchecked(exo10 e, int age) {
        if (age < 18) {
            throw new RuntimeException();
        }
    }
}
\end{lstlisting}

\subsection{Comparaison}

\begin{table}[H]
\centering
\begin{tabularx}{0.9\textwidth}{C{3.5cm} X X}
\toprule
\textbf{\color{primarycolor}Critère} &
\textbf{\color{primarycolor}Exception (Checked)} &
\textbf{\color{primarycolor}RuntimeException (Unchecked)} \\
\midrule
Vérification & À la compilation & Seulement à l'exécution \\
\addlinespace[0.15cm]
Signature & Oblige \texttt{throws} & Aucun \texttt{throws} requis \\
\addlinespace[0.15cm]
Usage & Erreurs récupérables externes & Bugs logiques du développeur \\
\bottomrule
\end{tabularx}
\caption{Comparaison entre \texttt{Exception} et \texttt{RuntimeException}}
\end{table}
